% ============================================================
%  BAB 2 – LANDASAN TEORI
% ============================================================


\chapter*{TINJAUAN PUSTAKA}
\refstepcounter{chapter}
\addcontentsline{toc}{chapter}{TINJAUAN PUSTAKA}
\pdfbookmark[0]{TINJAUAN PUSTAKA}{tinjauan_pustaka}


\section{Penelitian Terkait}
Dalam beberapa tahun terakhir berbagai solusi berbasis teknologi telah banyak dikembangkan untuk membantu individu dengan \textit{Color Vision Deficiency} (CVD) untuk meningkatkan kemampuan dalam membedakan warna.
Penelitian yang dilakukan oleh \textcite{Elrefaei2018} melakukan perbandingan empat algoritma koreksi warna berbasis \textit{smartphone}, termasuk algoritma \textit{LMS Daltonization}, yang dievaluasi menggunakan software simulator Vischeck dan Coblis.
Namun, pengujian pada penelitian tersebut memiliki keterbatasan karena pemrosesan citra masih berbasis statis dengan metode unggah gambar secara manual dan tidak berjalan secara real-time.

Selain itu, penelitian lain oleh \textcite{covisance} mengembangkan aplikasi seluler Covisance yang menerapkan metode LMS Daltonization secara real-time dan mengintegrasikan uji Farnsworth D-15 langsung di dalam aplikasi.
Meskipun mampu menampilkan simulasi secara langsung, aplikasi tersebut masih menggunakan metode kalkulasi nilai warna secara manual, yang dapat membebani kinerja CPU dan sering menyebabkan frame-drop pada perangkat dengan resolusi kamera tinggi.

Oleh karena itu, penelitian ini berfokus pada efisiensi komputasi visual secara realtime dengan cara memindahkan beban komputasi dari CPU ke GPU untuk menghindari proses sekuensial yang berat.
Dengan memanfaatkan \textit{library} CameraX dan \textit{ColorMatrix}, Dalsilens dapat mempertahankan stabilitas koreksi warna pada pemrosean \textit{real-time}, sehingga dapat mengurangi masalah penurunan bingkai (\textit{frame drop}) dan latensi yang sering terjadi pada pemrosesan citra real-time.

% ============================================================
\section{Konsep Dasar Penglihatan dan Warna}

\subsection{Color Vision Deficiency (CVD) dan Dikromasi}
\textit{Color Vision Deficiency} (CVD) atau buta warna adalah sebuah kondisi kelainan genetik yang memengaruhi kemampuan seseorang dalam membedakan panjang gelombang cahaya tertentu seperti merah, hijau, dan biru.
Tipe paling umum dari kondisi ini adalah Dikromasi, yang terbagi menjadi \textit{Protanopia} (buta warna merah), \textit{Deuteranopia} (buta warna hijau), dan \textit{Tritanopia} (buta warna biru).

\subsection{Ruang Warna Visual (RGB, LMS, dan HSV)}
Umumnya pemrosesan citra digital menggunakan ruang warna RGB (\textit{Red, Green, Blue}).
Namun, sistem penglihatan manusia memproses spektrum cahaya melalui sel kerucut mata yang direpresentasikan dalam ruang warna LMS (\textit{Long, Medium, Short}) \parencite{8289076}.
Selain itu, terdapat ruang warna HSV (\textit{Hue, Saturation, Value}) yang efektif dalam pemisahan intensitas cahaya dari warna itu sendiri.
Ruang warna ini membagi elemen visual menjadi tipe jenis warna (H), intensitas warna (S), dan tingkat kecerahan (V), sehingga cocok digunakan untuk klasifikasi warna.

% ============================================================
\section{Pengolahan Citra Digital (Digital Image Processing)}

\subsection{Pemrosesan Citra Real-Time}
Pengolahan citra digital merupakan teknik untuk memproses dan memodifikasi gambar digital dengan mengubah nilai piksel melalui proses komputasi.
Pengolahan citra digital secara real-time mengharuskan sistem untuk terus menerus mengambil sample bingkai (\textit{frame}) dari perangkat keras kamera.
Nilai piksel dari bingkai yang ditangkap tersebut kemudian diekstrak untuk diproses dan dianalisis sebelum ditampilkan kembali kepada pengguna secara real-time.

\subsection{Algoritma LMS Daltonization}
\textit{LMS Daltonization} adalah sebuah teknik pemrosesan gambar untuk memanipulasi warna pada citra, yang dapat membantu individu dengan buta warna supaya dapat membedakan warna tertentu.
Secara garis besar metode ini dibagi menjadi beberapa tahap, tahap pertama yaitu mengonversi ruang warna RGB menjadi ruang warna LMS.
Tahap kedua

% ============================================================
\section{Teknologi dan Perangkat Pengembang}

\subsection{Sistem Operasi Android}
Jelaskan Android sebagai platform.

\subsection{Pustaka Grafis Color Matrix dan Akselerasi Perangkat Keras}
Jelaskan penggunaan Color Matrix dan hardware acceleration.

% ============================================================